\begin{proof}
\textbf{Trivial cases.}
\begin{itemize}
\item \textbf{If $a\le 0$.}  
Since $X\ge 0$ a.s., the event $\{X\ge a\}$ is the whole space $\Omega$, hence
\[
\mathbb{P}(X\ge a)=1.
\]
Because $a\le 0$ the right–hand side $\frac{\mathbb{E}[X]}{a}$ is non‑negative (indeed it is $+\infty$ or undefined), so the inequality holds vacuously. From now on we restrict to $a>0$.

\item \textbf{If $\mathbb{E}[X]=0$.}  
For a non‑negative random variable, $\mathbb{E}[X]=0$ forces $X=0$ $\mathbb{P}$‑a.s. Consequently $\mathbb{P}(X\ge a)=0$ for any $a>0$, and the inequality becomes $0\le 0$, i.e. it holds with equality.
\end{itemize}
Thus we may assume
\[
a>0,\qquad \mathbb{E}[X]>0 .
\]

\medskip
\textbf{Indicator of the event $\{X\ge a\}$.}
Define the indicator (characteristic) function
\[
I(\omega):=\mathbf 1_{\{X\ge a\}}(\omega)=
\begin{cases}
1, & \text{if } X(\omega)\ge a,\\[4pt]
0, & \text{if } X(\omega)< a .
\end{cases}
\]
The indicator is measurable, takes values only in $\{0,1\}$, and by the
\emph{Expectation of an indicator} theorem,
\begin{equation}
\mathbb{E}[I]=\mathbb{P}\bigl(X\ge a\bigr). \tag{1}
\end{equation}

\medskip
\textbf{Pointwise comparison.}
For every $\omega\in\Omega$ we have
\begin{equation}
X(\omega)\;\ge\;a\,I(\omega). \tag{2}
\end{equation}
Indeed,
\begin{itemize}
\item If $X(\omega)\ge a$, then $I(\omega)=1$ and (2) reads $X(\omega)\ge a$, true.
\item If $X(\omega)<a$, then $I(\omega)=0$ and (2) reads $X(\omega)\ge 0$, true because $X$ is non‑negative.
\end{itemize}
Thus (2) holds pointwise.

\medskip
\textbf{Take expectations (monotonicity of the integral).}
The expectation (Lebesgue integral) is monotone: from (2) we obtain
\[
\mathbb{E}[X]\;\ge\;\mathbb{E}[a\,I].
\]
Since $a$ is a deterministic constant, linearity gives
\begin{equation}
\mathbb{E}[a\,I]=a\,\mathbb{E}[I]. \tag{3}
\end{equation}
Combining the two displays,
\begin{equation}
\mathbb{E}[X]\;\ge\;a\,\mathbb{E}[I]. \tag{4}
\end{equation}

\medskip
\textbf{Replace $\mathbb{E}[I]$ by the probability.}
Substituting \eqref{1} into \eqref{4} yields
\[
\mathbb{E}[X]\;\ge\;a\,\mathbb{P}\bigl(X\ge a\bigr).
\]
Dividing by the positive constant $a$ we obtain the desired inequality
\[
\boxed{\;\mathbb{P}\bigl(X\ge a\bigr)\le\frac{\mathbb{E}[X]}{a}\;}.
\]

\medskip
\textbf{Equality case.}
Equality in the chain
\[
X(\omega)\ge a\,I(\omega)\quad\Longrightarrow\quad\mathbb{E}[X]=a\,\mathbb{E}[I]
\]
requires equality in the pointwise inequality \eqref{2} for $\mathbb{P}$‑almost every $\omega$.
Hence, $\mathbb{P}$‑a.s.,
\[
X(\omega)=a\,I(\omega).
\]
Because $I(\omega)\in\{0,1\}$, this means that $X$ takes only the two values $0$ and $a$ (the degenerate sub‑case $X\equiv0$ corresponds to $\mathbb{P}(X\ge a)=0$). In this situation
\[
\mathbb{E}[X]=a\,\mathbb{P}(X\ge a),
\]
so the bound is attained.

\medskip
\textbf{Conclusion.}
Under the sole hypotheses $X\ge 0$ a.s.\ and $a>0$, Markov’s inequality
\[
\mathbb{P}(X\ge a)\le\frac{\mathbb{E}[X]}{a}
\]
holds, with equality if and only if $X$ is almost surely supported on $\{0,a\}$ (including the trivial case $X\equiv0$). \qed
\end{proof}